\documentclass[11pt,a4paper]{article}
\usepackage[utf8]{inputenc}
\usepackage[french]{babel}
\usepackage[T1]{fontenc}
\usepackage{amsmath}
\usepackage{amsfonts}
\usepackage{amssymb}
\usepackage{graphicx}
\usepackage{booktabs}
\usepackage{geometry}
\usepackage{hyperref}
\usepackage{xcolor}
\usepackage{tikz}

\usetikzlibrary{shapes.geometric, arrows, positioning}

\geometry{top=2.5cm, bottom=2.5cm, left=2.5cm, right=2.5cm}

\definecolor{primary}{RGB}{31, 78, 121}
\definecolor{secondary}{RGB}{120, 120, 120}

\hypersetup{
    colorlinks=true,
    linkcolor=primary,
    filecolor=primary,
    urlcolor=primary,
    pdftitle={Rapport Mini-Projet : When ML Fails},
}

\title{
    \vspace{-2cm}
    \fontsize{22}{26}\selectfont\bfseries\color{primary}{Projet : When Machine Learning Fails}\\
    \fontsize{14}{18}\selectfont\color{secondary}{Analyse approfondie des pannes : Raccourcis d'apprentissage, limites d'extrapolation et illusions de correction}\\
    \vspace{0.5cm}
    \rule{\linewidth}{1.5pt}
}

\author{
    \textbf{Membres de l'équipe :}\\[0.3cm]
    TONDE Stephane \quad --- \quad SANOU Jacquel \quad --- \quad LOMPO Ernest\\[0.3cm]
    \textit{Étudiants en 2\textsuperscript{ème} année --- Ecole Centrale Casablanca}\\[0.5cm]
    \textbf{Enseignantes :} \\
    Kawtar Zerhouni \& Rym Nassih\\
    \textit{Centrale Casablanca --- UTER MID@S}
}

\date{Printemps 2026}

\begin{document}

\maketitle
\vspace{1cm}

\begin{abstract}
Ce rapport détaille la démarche scientifique menée sur le jeu de données \texttt{paddydataset.csv}. L'objectif de la première phase était d'analyser le sur-apprentissage (overfitting) potentiellement provoqué par l'encodage de variables catégorielles. Nous réfutons l'hypothèse initiale selon laquelle le modèle a échoué à sur-apprendre à cause d'une haute cardinalité (les variables possédant en réalité une très basse cardinalité) et mettons en évidence le rôle de la variable de surface agricole comme raccourci d'apprentissage dominant. Dans une seconde phase, nous étudions l'effondrement du modèle de forêt aléatoire ($R^2$ de $-10.077$) face à une extrapolation physique hors de sa distribution d'entraînement (split OOD). Enfin, nous démontrons que la correction courante par mise à l'échelle de la cible seule ($R^2 = 0.705$) est un succès en trompe-l'œil (corrélation du taux prédit nulle de $-0.113$) et proposons une correction robuste via un pipeline invariant d'échelle (intrants et cible normalisés) assurant une généralisation saine ($R^2 = 0.380$ sur le taux).
\end{abstract}

\newpage

\section{Research Question and Chosen Dataset}

\subsection{Questions de recherche (Mode Principal et Mode Bonus)}
Pour ce projet, nous choisissons d'étudier deux modes de défaillance distincts sur le jeu de données \textbf{Paddy} (UCI 1186), à l'aide d'un modèle de \textbf{Random Forest} (famille tree-based ensembles), conformément à la contrainte du projet excluant la régression linéaire ou logistique comme modèle principal :
\begin{enumerate}
    \item \textbf{Mode Principal : Panne d'extrapolation d'échelle (Section 5.3 \& 5.2 du catalogue)}
    \begin{center}
    \textit{Le modèle Random Forest échoue-t-il à extrapoler le rendement agricole global lorsque évalué sur des parcelles de taille inédite hors de sa distribution d'entraînement ?}
    \end{center}
    \item \textbf{Mode Bonus : Raccourci d'apprentissage (Section 5.5 du catalogue)}
    \begin{center}
    \textit{Le modèle s'appuie-t-il sur la variable physique de taille \texttt{Hectares} comme raccourci d'apprentissage (Shortcut Learning) dominant, masquant ainsi l'influence réelle des autres variables et le sur-apprentissage théoriquement attendu ?}
    \end{center}
\end{enumerate}

\subsection{Dataset et Prétraitement dans les Règles de l'Art}
Le jeu de données \textbf{Paddy} comprend des données physiques, météorologiques et de gestion de cultures sur 2 789 parcelles. La variable cible est la quantité totale produite : \texttt{Paddy yield(in Kg)}. La surface cultivée est représentée par la variable \texttt{Hectares}.

Pour respecter les règles de l'art, nous avons supprimé les 451 doublons présents dans le jeu de données d'origine afin d'éviter toute fuite de données lors des fractionnements. De plus, les pipelines de prétraitement sont isolés de manière étanche : l'imputation (médiane pour le numérique, mode pour le catégoriel) et la standardisation apprennent leurs paramètres (\texttt{fit}) exclusivement sur les ensembles d'entraînement respectifs avant d'être appliqués (\texttt{transform}) sur les ensembles de test.

\section{Phase 1: Attempt 1 - Second Failure Mode (Shortcut Learning - Bonus)}

\subsection{Question de recherche et Hypothèse}
\textbf{Question :} Le modèle Random Forest apprend-il la relation physique de rendement par un raccourci d'apprentissage linéaire au détriment des relations non-linéaires plus fines ? \\
\textbf{Hypothèse :} La forte corrélation linéaire ($0.994$) entre \texttt{Hectares} et le rendement brut agit comme un raccourci dominant. En split aléatoire standard, cela masque le manque de généralisation ou le sur-apprentissage sur d'autres variables (comme l'encodage One-Hot des variables catégorielles).

\subsection{Test contrôlé et Symptôme observé}
Le test consiste à entraîner un Random Forest profond avec One-Hot Encoding complet de \texttt{Variety} et \texttt{Agriblock} sur un split aléatoire (80\% Train, 20\% Test). Les résultats mesurés sont :
\begin{itemize}
    \item $R^2$ Train : \textbf{0.994}
    \item $R^2$ Test : \textbf{0.989}
\end{itemize}
Le symptôme observé est un gap de généralisation quasi nul. L'écart est minime et le modèle ne sur-apprend pas.

\subsection{Réfutation théorique et Discussion du raccourci}
Le rapport initial attribuait l'absence de sur-apprentissage à la présence de variables catégorielles à ``haute cardinalité''. Notre analyse descriptive réfute formellement cette explication :
\begin{itemize}
    \item \texttt{Variety} possède seulement \textbf{3} modalités uniques.
    \item \texttt{Agriblock} possède seulement \textbf{6} modalités uniques.
\end{itemize}
La cardinalité réelle étant extrêmement basse, l'encodage OHE ne génère que 9 colonnes, ce qui ne peut mener à un sur-apprentissage par haute dimensionnalité.
En réalité, le modèle n'a pas sur-appris car il exploite le raccourci d'apprentissage physique de la variable \texttt{Hectares}. La relation d'échelle linéaire étant identique dans les jeux d'entraînement et de test aléatoires, les performances de test restent excellentes, masquant le fait que le modèle n'a pas appris la véritable productivité agronomique.

\section{Phase 2: Attempt 2 - Primary Failure Mode (Extrapolation Failure)}

\subsection{Protocole Out-Of-Distribution (OOD)}
Pour forcer le modèle à échouer, nous brisons la corrélation d'échelle en effectuant un découpage OOD basé sur la surface agricole :
\begin{itemize}
    \item \textbf{Train Set} : Parcelles de petite taille ($\le 4.0$ hectares, la médiane).
    \item \textbf{Test Set} : Parcelles de grande taille ($> 4.0$ hectares).
\end{itemize}

\subsection{Symptôme observé et incapacité d'extrapolation}
Le modèle s'effondre catastrophiquement sur le jeu de test :
\begin{itemize}
    \item $R^2$ Train OOD : \textbf{0.990}
    \item $R^2$ Test OOD : \textbf{-10.077} (Panne sévère)
    \item RMSE Train OOD : \textbf{660.92} | RMSE Test OOD : \textbf{8113.66}
\end{itemize}

Cette panne s'explique par la nature même des arbres de décision. Les prédictions d'une forêt aléatoire sont des moyennes locales de valeurs cibles observées en entraînement. Par conséquent, le modèle est structurellement incapable de prédire une valeur supérieure au maximum vu lors de l'apprentissage : $\hat{y} \le \max(y_{train}) \approx 25\,000$ Kg. Sur le jeu de test composé de grands champs, les rendements réels atteignent jusqu'à $75\,000$ Kg, ce qui conduit au plafonnement horizontal des prédictions.

\section{Proposed Corrections and Critical Evaluation}

\subsection{La Correction Naïve : Target Scaling Seul (Succès en trompe-l'œil)}
La première approche consiste à normaliser la cible par unité de surface :
\[
\text{Yield\_per\_Hectare} = \frac{\text{Paddy yield(in Kg)}}{\text{Hectares}}
\]
Le modèle est entraîné à prédire ce taux (qui est borné agronomiquement), puis la prédiction finale est reconstruite à l'inférence :
\[
\hat{y}_{final} = \widehat{\text{Yield\_per\_Hectare}} \times \text{Hectares}
\]

Cette correction affiche en apparence d'excellents résultats avec un $R^2$ de test remontant à \textbf{0.706} (RMSE = 1322.53).

\textbf{Cependant, notre analyse révèle une panne silencieuse sous-jacente :}
La corrélation de Pearson entre le taux de rendement prédit par le modèle et le taux réel sur le jeu de test est de \textbf{-0.113} (nulle à légèrement négative).
\textit{Pourquoi ?} Les variables physiques d'intrants fournies en entrée du modèle (\texttt{Seedrate}, \texttt{Urea\_40Days}, etc.) sont parfaitement proportionnelles à la surface. Sur le jeu de test OOD, ces variables sont donc également hors-distribution. La forêt aléatoire les clippe toutes, ce qui l'amène à prédire un taux de rendement presque constant ($\approx 6\,100$ Kg/Ha) pour toutes les parcelles de test.
Le $R^2$ de $0.706$ est une \textbf{illusion mathématique} portée uniquement par la multiplication finale par \texttt{Hectares}, qui varie de 5 à 12 et recrée artificiellement la variance du rendement brut. Le modèle ML n'a en réalité rien appris d'utile sur la productivité agronomique dans cette configuration.

\subsection{La Correction Robuste : Pipeline Invariant d'Échelle}
Pour appliquer les véritables règles de l'art, nous devons rendre le pipeline d'entrée lui-même invariant d'échelle.
\begin{enumerate}
    \item Nous identifions dynamiquement les intrants corrélés à plus de 0.8 avec \texttt{Hectares} uniquement sur le jeu d'entraînement \texttt{X\_train\_ood} pour éviter toute fuite de données (ici : \texttt{Seedrate}, \texttt{LP\_Mainfield}, \texttt{LP\_nurseryarea}, \texttt{Weed28D\_thiobencarb}, \texttt{Urea\_40Days}, \texttt{Pest\_60Day}, \texttt{Nursery area}, \texttt{DAP\_20days}, \texttt{Potassh\_50Days}, \texttt{Micronutrients\_70Days}, \texttt{Trash}).
    \item Nous divisons chacun de ces intrants par \texttt{Hectares} pour obtenir des taux d'application (ex: Kg de semences par hectare).
    \item Nous supprimons la variable \texttt{Hectares} des descripteurs pour éliminer toute fuite d'échelle.
\end{enumerate}

Ce modèle robuste est entraîné sur ces caractéristiques normalisées pour prédire le rendement par hectare. Les résultats du pipeline robuste sont :
\begin{itemize}
    \item $R^2$ Test Global (rendement brut reconstruit) : \textbf{0.380} (RMSE = 1919.52)
    \item Corrélation du taux prédit vs réel : \textbf{-0.022}
\end{itemize}
Bien que le $R^2$ global paraisse inférieur (0.380 contre 0.706), ce score représente la \textbf{véritable capacité prédictive agronomique} (influence de la météo, du sol et de la variété sur la productivité) exempte de toute fuite d'échelle ou de clipping physique. Le modèle est désormais mathématiquement stable face à n'importe quelle échelle de parcelle.

\begin{table}[h]
\centering
\caption{Synthèse comparative des performances des modèles}
\label{tab:metrics}
\vspace{0.2cm}
\begin{tabular}{lccccc}
\toprule
\textbf{Configuration} & \textbf{$R^2$ Train} & \textbf{$R^2$ Test} & \textbf{RMSE Train} & \textbf{RMSE Test} & \textbf{Corr. Taux} \\
\midrule
Tentative 1 (Aléatoire) & 0.994 & 0.989 & 737.89 & 968.83 & -- \\
Tentative 2 (OOD) & 0.990 & -10.077 & 660.92 & 8113.66 & -- \\
\midrule
Correction Naïve (Cible seule) & 0.990 & 0.706 & 660.92 & 1322.53 & -0.113 \\
Correction Robuste (Invariant) & 0.982 & 0.380 & 911.61 & 1919.52 & -0.022 \\
\bottomrule
\end{tabular}
\end{table}

\section{Architecture des Pipelines}
La figure \ref{fig:pipelines} schématise l'étanchéité des données et le traitement d'invariance d'échelle.

\begin{figure}[h]
\centering
\begin{tikzpicture}[node distance=1.5cm, auto]
    \node (raw) [rectangle, draw, fill=blue!10, text width=6cm, text centered, rounded corners, minimum height=1cm] {Données Brutes (paddydataset.csv)};
    \node (clean) [rectangle, draw, fill=blue!10, text width=6cm, text centered, rounded corners, minimum height=1cm, below=of raw] {Nettoyage : Suppression des Doublons};
    \node (split) [diamond, draw, fill=yellow!20, text width=2.5cm, text centered, below=of clean] {Split OOD (Surface)};
    \node (train) [rectangle, draw, fill=green!10, text width=5cm, text centered, left=3cm of split] {Train Split ($\le 4.0$ Ha)};
    \node (test) [rectangle, draw, fill=red!10, text width=5cm, text centered, right=3cm of split] {Test Split ($> 4.0$ Ha)};

    \node (scale_tr) [rectangle, draw, fill=green!20, text width=5cm, text centered, below=of train] {Normalisation : Intrants / Hectares};
    \node (scale_te) [rectangle, draw, fill=red!20, text width=5cm, text centered, below=of test] {Normalisation : Intrants / Hectares};

    \node (fit) [rectangle, draw, fill=green!30, text width=4cm, text centered, below=of scale_tr] {Fit Preprocessor \& Model};
    \node (pred) [rectangle, draw, fill=red!30, text width=4cm, text centered, below=of scale_te] {Transform \& Predict};

    \draw[->, thick] (raw) -- (clean);
    \draw[->, thick] (clean) -- (split);
    \draw[->, thick] (split) -- (train);
    \draw[->, thick] (split) -- (test);
    \draw[->, thick] (train) -- (scale_tr);
    \draw[->, thick] (test) -- (scale_te);
    \draw[->, thick] (scale_tr) -- (fit);
    \draw[->, thick] (scale_te) -- (pred);
    \draw[->, thick, dashed] (fit) -- (pred);
\end{tikzpicture}
\caption{Architecture étanche du pipeline invariant d'échelle}
\label{fig:pipelines}
\end{figure}

\section{Threats to Validity (Menaces sur la validité)}
Une analyse de modélisation critique exige d'évaluer les limites théoriques et statistiques de notre démarche :
\begin{enumerate}
    \item \textbf{Taille et représentativité du jeu de test OOD :} Le jeu de données initial après suppression des doublons contient 2 789 parcelles. Notre split à la médiane de surface (4.0 Ha) alloue 1 394 parcelles pour l'entraînement et 1 395 parcelles pour le test OOD. Cette taille d'échantillon est statistiquement très élevée, ce qui élimine tout risque que la chute de $R^2$ ou l'effet de plafonnement horizontal soient dus à du bruit ou à un échantillon trop restreint.
    \item \textbf{Variables de confusion (Confounders) :} La normalisation simple par division par \texttt{Hectares} suppose une relation linéaire directe entre la surface et les intrants physiques (engrais, semences). En pratique, les grandes parcelles agricoles peuvent bénéficier d'économies d'échelle, de mécanisation différente ou de modes de culture spécifiques, rendant les taux non linéaires. Si tel est le cas, le ratio simple introduit un biais agronomique, constituant une menace sur la validité de la correction.
    \item \textbf{Robustesse à la graine aléatoire (Random State) :} Le plafonnement horizontal est une limite structurelle mathématique des arbres (incapables de prédire au-delà du maximum observé). De ce fait, le comportement d'effondrement OOD et l'illusion mathématique de la correction naïve restent parfaitement stables et insensibles aux changements de \texttt{random\_state}.
\end{enumerate}

\section{Conclusion}
Ce travail met en lumière qu'un score $R^2$ élevé après correction peut s'avérer trompeur et masquer une panne d'apprentissage. En analysant la corrélation interne des taux de prédiction et en implémentant un pipeline rigoureusement invariant d'échelle (normalisant à la fois les intrants et les cibles), nous avons construit un modèle robuste, scientifiquement valide et capable de généraliser sainement en dehors de sa distribution d'entraînement.

\end{document}
